\section{Syntax, semantics, and proof theory}
\label{sec:levels-03}

\begin{definition}[Syntax]
The \emph{syntax} of a formal language consists of its alphabet, grammar, and
operations on expressions defined without reference to interpretation.
\end{definition}

\begin{definition}[Proof theory]
The \emph{proof theory} of a calculus studies derivations and transformations
of derivations by syntactic methods.
\end{definition}

Syntax and proof theory are related but distinct.  Syntax determines which
formulas exist.  A calculus determines which derivations exist.  Semantics
determines when formulas are satisfied.  Metatheorems connect these domains.

\begin{definition}[Semantic consequence]
Let \(\Gamma\) be a set of formulas and let \(\varphi\) be a formula in a
language with a specified semantics.  Then \(\varphi\) is a \emph{semantic
consequence} of \(\Gamma\), written
\[
  \Gamma\entails\varphi,
\]
iff, for every complete tuple of semantic data permitted by that semantics,
satisfaction of every member of \(\Gamma\) implies satisfaction of
\(\varphi\).  The permitted data may be a propositional valuation, a
first-order structure with a variable assignment, a modal structure with a
world, or another tuple explicitly declared by the active semantics.
\end{definition}

\begin{definition}[Derivability]
Let \(\mathcal C\) be a specified calculus.  A formula \(\varphi\) is
\emph{derivable} from a set of formulas \(\Gamma\), written
\[
  \Gamma\derives_{\mathcal C}\varphi,
\]
iff there exists a derivation in \(\mathcal C\) whose conclusion is
\(\varphi\) and whose undischarged assumptions are members of \(\Gamma\).
\end{definition}

\begin{definition}[Validity]
A formula \(\varphi\) is \emph{valid} for a specified semantics iff
\(\varnothing\entails\varphi\).  We then write \(\entails\varphi\).
\end{definition}

\begin{definition}[Theorem of a calculus]
A formula \(\varphi\) is a \emph{theorem} of a specified calculus
\(\mathcal C\) iff \(\varnothing\derives_{\mathcal C}\varphi\).  We then
write \(\derives_{\mathcal C}\varphi\).
\end{definition}

\begin{definition}[Soundness]
A calculus \(\mathcal C\) is \emph{sound} for a specified semantics iff, for
every set of formulas \(\Gamma\) and formula \(\varphi\),
\[
  \Gamma\derives_{\mathcal C}\varphi
  \quad\Longrightarrow\quad
  \Gamma\entails\varphi.
\]
\end{definition}

\begin{definition}[Completeness]
A calculus \(\mathcal C\) is \emph{complete} for a specified semantics iff,
for every set of formulas \(\Gamma\) and formula \(\varphi\),
\[
  \Gamma\entails\varphi
  \quad\Longrightarrow\quad
  \Gamma\derives_{\mathcal C}\varphi.
\]
\end{definition}

\begin{warning}[Direction check]
Soundness begins with a derivation.  Completeness begins with semantic
consequence.  Writing the two displayed implications before beginning a proof
prevents reversal of their roles.
\end{warning}

